\VerdierPhysicalPageAnchor{0026}%
donnée et changement de signe de la différentielle correspondante). Les
extensions naturelles des foncteurs multi-additifs aux complexes ne commutent
pas nécessairement aux translations, mais y commutent \og à isomorphisme
près\fg{}, isomorphismes déterminés, pratiquement, au signe près. Il se trouve
qu'on ne peut choisir ces isomorphismes de manière cohérente pour la
composition des \og foncteurs translation\fg{}, et qu'il s'introduit alors un
signe. L'étude de ce signe fait l'objet du paragraphe 1 (voir aussi
l'appendice). Au paragraphe 2, nous définissons les notions classiques de
complexes multiples, complexes simples, foncteur \og complexe simple
associé\fg{}, homotopie, et au paragraphe 3, nous abordons l'étude du cône
d'un morphisme de complexes (\og mapping cone\fg{} ou \og mapping
cylinder\fg{}), construction essentielle pour la suite.

\VerdierSemanticAnchor{intro-6-plan-derived-category}
La construction de la catégorie $\mathbf D(\mathcal A)$, catégorie dérivée de
la catégorie abélienne $\mathcal A$, se fait en deux étapes. Dans la première
étape on introduit la catégorie $\mathbf K(\mathcal A)$ des complexes de
$\mathcal A$ à homotopie près. Puis on inverse formellement dans
$\mathbf K(\mathcal A)$ les quasi-isomorphismes pour obtenir la catégorie
$\mathbf D(\mathcal A)$. Les catégories $\mathbf K(\mathcal A)$ et
$\mathbf D(\mathcal A)$ ne sont pas nécessairement abéliennes; mais elles
sont munies d'une structure supplémentaire, consistant en la donnée d'une
famille de diagrammes définis à partir de la construction du cône :
\emph{les triangles distingués}. Les triangles distingués jouent, pour ces
catégories, le rôle des suites exactes des catégories abéliennes.

\VerdierSemanticAnchor{intro-6-plan-chapter-two}
Les catégories additives munies de cette structure supplémentaire, mise en
évidence et étudiée par Puppe [2], sont appelées
\emph{catégories triangulées} et sont étudiées de notre point de vue au
chapitre II. Le premier paragraphe est consacré à l'exposé de résultats
élémentaires sur les catégories triangulées. Au paragraphe 2, nous montrons,
dans le cadre général des catégories triangulées, que le problème consistant
à inverser formellement les quasi-isomorphismes de la catégorie triangulée
$\mathbf K(\mathcal A)$ se résout simplement par un \og calcul de
fractions\fg{} [16]. Au paragraphe 3, nous démontrons un théorème dû à Freyd
[17] : une catégorie triangulée $\mathcal D$ se plonge de manière universelle
dans une catégorie abélienne $\mathrm A(\mathcal D)$. Cette catégorie
$\mathrm A(\mathcal D)$ possède une famille $E$ d'objets à la fois injectifs
et projectifs, telle que tout objet de $\mathrm A(\mathcal D)$ soit isomorphe
à un sous-objet d'un objet de $E$ et à un objet quotient d'un objet de $E$.
La sous-catégorie pleine de $\mathrm A(\mathcal D)$ définie par $E$ est
isomorphe à la catégorie $\mathcal D$. Enfin au paragraphe 4, nous étudions
les \og objets spectraux\fg{} à valeurs dans les catégories triangulées ou
les catégories abéliennes. Nous ne faisons ici qu'adapter les définitions de
[1].

\VerdierSemanticAnchor{intro-6-plan-chapter-three}
Le chapitre III est consacré à l'étude des catégories dérivées proprement
dites. Lorsque $X$ et $Y$ sont deux objets de $\mathbf D(\mathcal A)$,
catégorie dérivée de la catégorie abélienne $\mathcal A$, les morphismes de
degré $n$ de $X$ dans $Y$ s'interprètent, lorsque la catégorie $\mathcal A$
possède suffisamment d'objets injectifs et lorsque $Y$ est un complexe dont
les composants sont nuls en degré petit, comme les

